\section*{Practical Appendices}
\addcontentsline{toc}{section}{Practical Appendices}
\label{sec:practical_appendices_adjacencies}

\subsection*{Appendix A: Adjacency Screening Card}
\addcontentsline{toc}{subsection}{Appendix A: Adjacency Screening Card}
\label{sec:appendix_a_adjacency_screening_card}

\begin{tcolorbox}[colback=white,colframe=black!75,arc=1mm,title=\textbf{Adjacency Screening Card: Initial Feasibility Filter}]
\small
\noindent\textbf{Opportunity / Owner / Date:} \\
\textbf{Target Buyer, Unmet Job, and Purchase Trigger:} \\
\textbf{What Specifically Changes from Our Existing Offering:} \\
\textbf{Closest Credible Substitute (including non-consumption):} \\
\textbf{Empirical Evidence of Paid Demand (not survey intent):} \\
\textbf{Transferred Core Asset and What Cost/Risk It Reduces:} \\
\textbf{New Capabilities and Specialist Staffing Required:} \\
\textbf{Incremental Economics, Founder Hours, and Cash Timing:} \\
\textbf{Operational, Ethical, and Legal Hard Gates:} \\
\textbf{What Durable Institutional Asset Could Accumulate:} \\
\textbf{Best Alternative Use of Our Capital and Founder Time:} \\
\textbf{Smallest Reversible Experiment to Resolve Top Risk:} \\
\textbf{Predefined Pass, Fail, and Revisit Triggers:}
\end{tcolorbox}

\subsection*{Appendix B: The 30-Minute Strategic Conversation}
\addcontentsline{toc}{subsection}{Appendix B: The 30-Minute Strategic Conversation}
\label{sec:appendix_b_thirty_minute_conversation}

Ask the proposer to speak 1st for 5 minutes without slides. Then ask:
\begin{enumerate}
    \item What real customer decision produced the opportunity?
    \item What does the customer use today?
    \item Why is this firm better positioned than an unrelated specialist?
    \item Which current work would stop if we pursue this?
    \item What ongoing operational duty is created?
    \item What cash is due before collection?
    \item Who would operate it if the founder were absent?
    \item Which specific evidence would make the proposer abandon it?
\end{enumerate}
Close by assigning 1 concrete test, not an unbounded discovery initiative.

\subsection*{Appendix C: A Reversible Launch Contract Checklist}
\addcontentsline{toc}{subsection}{Appendix C: A Reversible Launch Contract Checklist}
\label{sec:appendix_c_reversible_launch_contract}

For a small adjacency pilot document:
\begin{itemize}
    \item Scope boundary and explicit out-of-scope exclusions;
    \item Objective acceptance criteria and sign-off authority;
    \item Payer identity and milestone-based payment schedules;
    \item Clear cancellation notice and termination liabilities;
    \item Defined post-delivery support term and SLA boundaries;
    \item Explicit ownership of foreground IP, background IP, and learner data;
    \item Confidentiality protections and non-solicitation of key talent;
    \item Scope-change governance and change-order pricing;
    \item Permitted use of institutional marks, logos, and case references;
    \item Clean exit and data-return protocols upon completion.
\end{itemize}
Avoid making commitments through sales decks that the written contract contradicts. Obtain professional review for consequential terms.

\subsection*{Appendix D: Scope-Economy Experiment}
\addcontentsline{toc}{subsection}{Appendix D: Scope-Economy Experiment}
\label{sec:appendix_d_scope_economy_experiment}

Choose 1 shared asset: instructor calibration, an assessment platform, reusable code, distribution, or supplier procurement. Compare total quality-adjusted cost when Program A and Program B use it jointly versus separately:
\begin{equation}
\label{eq:scope_economy_test}
\Delta C_{\text{scope}} = [C(A,0) + C(0,B)] - C(A,B)
\end{equation}
Include coordinating staff, delays, rework, and opportunity cost. If the shared asset creates no measurable saving or better outcome ($\Delta C_{\text{scope}} \le 0$), do not justify business expansion by invoking economies of scope.

\subsection*{Appendix E: Portfolio Red-Team Questions}
\addcontentsline{toc}{subsection}{Appendix E: Portfolio Red-Team Questions}
\label{sec:appendix_e_portfolio_red_team}

Stress-test the multi-business portfolio against 7 adversarial scenarios:
\begin{enumerate}
    \item What happens if the largest buyer leaves tomorrow?
    \item What happens if the 2 main products depend on the same underlying technology cycle?
    \item What happens if the founder is completely unavailable for 3 months?
    \item What happens if institutional cash receipts slip by 90 days?
    \item What happens if a promising research project fails after spending its entire allocation?
    \item What happens if a flagship institutional partner requests exclusivity?
    \item What happens if a new product succeeds at volume but lowers the institution's academic standards?
\end{enumerate}
Write concrete mitigation answers before these scenarios occur.

\subsection*{Appendix F: A 60-Day Adjacency Discovery Sprint}
\addcontentsline{toc}{subsection}{Appendix F: A 60-Day Adjacency Discovery Sprint}
\label{sec:appendix_f_sixty_day_discovery_sprint}

\begin{itemize}
    \item \textbf{Days 1-10:} Map alternatives, customer context, existing capability, and hard non-compensatory gates.
    \item \textbf{Days 11-20:} Conduct buyer and stakeholder discovery; secure concrete evidence of budgets and purchase authority.
    \item \textbf{Days 21-30:} Design a paid, reversible minimum offer and evaluate delivery capacity.
    \item \textbf{Days 31-45:} Sell or formally test the offer; collect process and pricing elasticity evidence.
    \item \textbf{Days 46-55:} Assess outcome quality, economic contribution, founder attention, and transferable assets.
    \item \textbf{Days 56-60:} Produce an expansion, hold, redesign, or stop memo.
\end{itemize}
Adjust for industry sales cycles: a 60-day interval may be too short to complete formal university procurement.

\subsection*{Appendix G: The Integration Break-Even Tool}
\addcontentsline{toc}{subsection}{Appendix G: The Integration Break-Even Tool}
\label{sec:appendix_g_integration_breakeven_tool}

Let external unit price be $P_e$, internal annual fixed overhead be $F$, and internal incremental cost per unit be $V_i$. Assuming comparable output quality and available capacity:
\begin{equation}
\label{eq:appendix_breakeven_formula}
Q^* = \frac{F}{P_e - V_i}, \quad \text{provided } P_e > V_i
\end{equation}
Stress-test below, at, and above $Q^*$. Add recruitment, management, idle capacity, variance, quality, and capital costs. Consider a hybrid solution when volume is intermittent. The formula is an entry point for analysis, not a hiring authorization.

\subsection*{Appendix H: The Strategic Parking Lot}
\addcontentsline{toc}{subsection}{Appendix H: The Strategic Parking Lot}
\label{sec:appendix_h_strategic_parking_lot}

Maintain a dated register of attractive rejected opportunities:
\begin{itemize}
    \item Specific rationale for refusal today;
    \item Evidential or readiness gap required to reconsider;
    \item Scarce founder and organizational resources required;
    \item Specific external market trigger for review;
    \item Alignment with long-term institutional mission.
\end{itemize}
Review quarterly. This limits recency bias and prevents founders from treating refusal as loss of ambition. Remove ideas that no longer make sense.

\subsection*{Appendix I: Bayes-Specific Adjacency Ledger}
\addcontentsline{toc}{subsection}{Appendix I: Bayes-Specific Adjacency Ledger}
\label{sec:appendix_i_bayes_adjacency_ledger}

Track each potential route separately: an assessed learner studio, a 2nd technical module, a college elective, faculty development, enterprise AI training, workflow consulting, a cross-border learner pilot, assessment services, research tools, and eventual physical infrastructure.

For each line maintain: target buyer, proven demand, delivery quality, founder hours, cash timing, IP rights, institutional asset transfer, and next smallest experiment. Do not report all as parts of 1 proven business until empirical evidence supports it.

\subsection*{Appendix J: 8 Mistakes to Avoid}
\addcontentsline{toc}{subsection}{Appendix J: 8 Mistakes to Avoid}
\label{sec:appendix_j_eight_mistakes_to_avoid}

\begin{enumerate}
    \item Confusing total market size with achievable economic entry.
    \item Counting shared staff as economies of scope despite severe bottlenecks.
    \item Funding an unvalidated new business line from tuition delivery cash reserves.
    \item Claiming every AI-adjacent consulting project is strategically related.
    \item Assuming revenue synergies without verified buyer purchasing behavior.
    \item Acquiring a company without dedicated integration management capacity.
    \item Believing geographic proximity makes regulatory and tax differences irrelevant.
    \item Assuming an experiment preserves options when it creates binding service promises.
\end{enumerate}

\section*{Research Foundations and Further Reading}
\addcontentsline{toc}{section}{Research Foundations and Further Reading}
\label{sec:research_foundations_adjacencies}

The distinction between source-derived frameworks and this chapter's managerial syntheses matters. The following works support major theoretical lenses; no individual paper validates the Bayes-specific scenarios or worksheets.

\begin{enumerate}
    \item \textbf{H. Igor Ansoff (1957), ``Strategies for Diversification,''} \textit{Harvard Business Review}. Classical product-market growth framework.
    \item \textbf{C. K. Prahalad and Gary Hamel (1990), ``The Core Competence of the Corporation,''} \textit{Harvard Business Review}. Competence-based approaches to multi-business strategy.
    \item \textbf{Michael E. Porter (1996), ``What Is Strategy?''} \textit{Harvard Business Review}. Trade-offs and interlocking activity systems.
    \item \textbf{David J. Teece (1986), ``Profiting from Technological Innovation,''} \textit{Research Policy}. Complementary assets, integration, collaboration, and value capture.
    \item \textbf{Ronald H. Coase (1937), ``The Nature of the Firm,''} \textit{Economica}. Market coordination and the rationale for firm boundaries.
    \item \textbf{Oliver E. Williamson (1985), \textit{The Economic Institutions of Capitalism}.} Transaction costs, asset specificity, and governance choices.
    \item \textbf{Rita Gunther McGrath and Ian C. MacMillan (1995), ``Discovery-Driven Planning,''} \textit{Harvard Business Review}. Assumption-driven staged investment and milestone gating.
    \item \textbf{Edith T. Penrose (1959), \textit{The Theory of the Growth of the Firm}.} Managerial resources, unused administrative services, and limits to expansion rate.
    \item \textbf{Richard P. Rumelt (1974), \textit{Strategy, Structure, and Economic Performance}.} Related and unrelated diversification; strategic and organizational fit.
    \item \textbf{Jay B. Barney (1991), ``Firm Resources and Sustained Competitive Advantage,''} \textit{Journal of Management}. Resource advantages and imitation constraints.
    \item \textbf{Jeffrey H. Dyer and Harbir Singh (1998), ``The Relational View,''} \textit{Academy of Management Review}. Interorganizational rents and complementary resources.
    \item \textbf{Ron Adner (2017), ``Ecosystem as Structure,''} \textit{Journal of Management}. Interdependencies required for a focal value proposition.
\end{enumerate}

\vspace{1em}
\begin{tcolorbox}[colback=black!5,colframe=black!75,arc=1mm,title=\textbf{The Final Principle}]
\textbf{The right adjacency should make tomorrow's institution more capable, not merely make today's business larger.}

\vspace{0.5em}
A strategically disciplined founder does not confuse nearby ideas with transferable capabilities, multiple revenues with diversification, ownership with control, or experimentation with freedom from obligations. The organization earns the right to expand by repeatedly making its current promises true, learning what it can reuse, and choosing new commitments that strengthen the system rather than scattering it.

That question leads directly to Chapter 30: once value depends on capabilities outside the company's boundary, \textbf{how should the enterprise design partnerships and ecosystems in which all necessary participants have reason to make the promise real?}
\end{tcolorbox}
